\documentclass[12pt,a4paper]{article}

\setlength{\topmargin}{0.0in}
\setlength{\oddsidemargin}{0.33in}
\setlength{\textheight}{9.0in}
\setlength{\textwidth}{6.0in}
\renewcommand{\baselinestretch}{1.25}



\usepackage{hyperref}
% \usepackage[utf8]{inputenc}
% \usepackage[T1]{fontenc}
\usepackage{float}
\usepackage{enumitem}
\usepackage{listings}
\usepackage{amsmath}
\usepackage{amssymb}
\usepackage{amsthm}
\usepackage{tikz}
\usetikzlibrary{decorations.pathreplacing}
\usepackage{amsfonts}
\usepackage{graphicx}
\usepackage{longtable}
\usepackage{booktabs}
\usepackage{mathtools}
\usepackage{subcaption}
\usepackage{caption}
\newcommand{\dv}[2]{\frac{\mathrm{d} #1}{\mathrm{d} #2}}
\usepackage{caption}
\captionsetup[figure]{font=small} 
% \captionsetup[table]{font=small}
\usepackage[capitalize,nameinlink]{cleveref}
\usepackage{mathtools}
\newcommand{\defeq}{\vcentcolon=}
\newcommand{\eqdef}{=\vcentcolon}
\usepackage{amsmath}
\usetikzlibrary{arrows.meta}

\definecolor{meshOrange}{RGB}{238,145,0}
\definecolor{cellGreen}{RGB}{135,220,20}
\definecolor{eigenPurple}{RGB}{155,0,190}
\definecolor{axisGray}{RGB}{75,85,85}

% Define environments
% \newtheorem{theorem}{Theorem}
% \newtheorem{definition}{Definition}
% \newtheorem{lemma}{Lemma}
% \newtheorem{corollary}{Corollary}


\newtheoremstyle{boldthm}
  {3pt}        % space above
  {3pt}        % space below
  {\itshape}   % body font
  {}           % indent
  {\bfseries}  % theorem head font
  {.}          % punctuation
  {.5em}       % space after theorem head
  {\thmname{#1}\thmnumber{ #2}\thmnote{ \textbf{(#3)}}}

\newtheoremstyle{bolddefn}
  {3pt}
  {3pt}
  {\itshape} 
  {}
  {\bfseries}
  {.}
  {.5em}
  {\thmname{#1}\thmnumber{ #2}\thmnote{ \textbf{(#3)}}}


\theoremstyle{boldthm}
\newtheorem{theorem}{Theorem}[section]
\newtheorem{corollary}[theorem]{Corollary}
\newtheorem{lemma}[theorem]{Lemma}
\newtheorem{proposition}[theorem]{Proposition}

\theoremstyle{bolddefn}
\newtheorem{definition}[theorem]{Definition}
\newtheorem{example}[theorem]{Example}
\newtheorem{remark}[theorem]{Remark}



\title{Speeding up Outer Loop for Self-Adjoint Problems}
\author{}
\date{}

\begin{document}

\maketitle

\section{Matt's Folded Operator}

Let $z = x + iy.$ Matt's folded operator is $F(z) \defeq (A-zI)^*(A-zI).$ Assume that $A^*=A.$ Then,
$F(z) = (A-\overline{z}I)(A-zI).$

Next, plug in $z$ and expand
\begin{align}
    F(z) &= (A-xI + iyI)(A-xI-iyI)\\
    &=(A-zI)^2 + y^2 I.
\end{align}

Furthermore, for the injection modulus, we have
\begin{align}
    \|(A-zI) u \|^2 &= \| (A-zI) u - iyu \|^2 \\
    &= \| (A-xI)u \|^2 + y^2 \|u\|^2 + 2 \Re\{ (A-xI)u, -iyu ) \} \\
    &= \| (A-xI)u \|^2 + y^2 \|u\|^2.
\end{align}
Taking the infinum over unit vectors, we obtain
\begin{equation}
    \gamma(z, A)^2 = \gamma(x, A)^2 + y^2. 
\end{equation}

\section{FEM Eigenproblem}

Next, we consider $B_z \defeq (A-zM)^*M^{-1}(A-zM) = (A-\overline{z}M)M^{-1}(A-zM).$ Expanding, we have 
\begin{align}
    B_z &= AM^{-1}A - (z + \overline z) A + |z|^2 M \\
    &= AM^{-1} A - 2xA + (x^2 + y^2) M.
\end{align}
Notice that 
\begin{equation}
    B_x = (A-xM)M^{-1}(A-xM) = AM^{-1}A - 2xA + x^2 M.
\end{equation}
Thus, 
\begin{equation}
    B_z = B_x + y^2 M.
\end{equation}
Now, define the sesquilinear forms 
\begin{align}
    a_z(u_h, v_h) &= \langle B_z u_h, v_h \rangle \\
    &= \langle B_x u_h, v_h \rangle + y^2 \langle M u_h, v_h \rangle  \\
    &= a_x(u_h, v_h) + y^2 m(u_h, v_h).
\end{align}
Thus, the eigenproblem 
\begin{equation}
    a_z(u_h, v_h) = \lambda_h(z) m(u_h, v_h),
\end{equation}
can be written as 
\begin{equation}
    a_x(u_h, v_h) = (\lambda_h(z) - y^2) m(u_h, v_h) = \lambda_h(x) m(u_h, v_h).
\end{equation}
In other words, for each distinct $x,$ we can solve $a_x(u_h, v_h) = \lambda_h(x) m(u, v)$ and find $\lambda_h(z)$ by adding $y^2.$ Furthermore, the eigenfunctions carry over as well. That is, $E_h(z) = E_h(x),$ where $E_h(\cdot)$ denotes the eigenspace for $\lambda_h(\cdot).$

With this, we can conclude that for $z$ values that share the same real part, Matt's bound on the distance to the spectrum depends only on a shift by the square of the imaginary part:
\begin{equation}
    \Phi_n(z, A) = \sqrt{\lambda_h(z)} = \sqrt{\lambda_h(x) + y^2}.
\end{equation}


\section{DWR Error Estimate}

Notice that the imaginary part of $z$ does not affect the primal residual:
\begin{align}
    \rho_z(v) &\defeq a_z(u_h, v) - \lambda_h(z) m(u_h, v) \\
    &= a_x(u_h, v) + y^2 m(u_h, v) - \lambda_h(z) m(u_h, v) \\
    &= a_x(u_h, v) - (\lambda_h(z) - y^2) m(u_h, v) \\
    &= a_x(u_h, v) - \lambda_h(x) m(u_h, v) \\
    &= \rho_x(v).  
\end{align}
Since the residual only depends on $x,$ and the eigenfunctions only depend on $x,$ the DWR error estimate depends only on $x$:
\begin{align}
    \lambda(z) - \lambda_h(z) &\approx \frac{\rho_z(u_h^+(z) - u_h(z))}{1 - 1/2 \| u_h^+(z) - u_h(z) \|_m^2 }\\
    &= \frac{\rho_z(u_h^+(x) - u_h(x))}{1 - 1/2 \| u_h^+(x) - u_h(x) \|_m^2 }.
\end{align}


\section{What does this mean?}

If correct, this means we can drastically speed up our computations. For each $x,$ we do the expensive adaptive eigensolves. Then, for each `column' (every $z$ that shares $x$) in the complex plane, we compute the needed $\Phi_n(z, A)$ using the shift by $y^2.$ The error estimate computed at $x$ carries over to every other point in that column.

\end{document}